\QHeading{第1题　口算。}
\par\noindent


\begin{oralgrid}
\OralItem{1}{16\times6}{96}
\OralItem{2}{900\times0}{0}
\OralItem{3}{420\times20}{8400}
\OralItem{4}{20\times300}{6000}
\OralItem{5}{32+140}{172}
\OralItem{6}{450-21}{429}
\OralItem{7}{90\times70}{6300}
\OralItem{8}{95\div5}{19}
\OralItem{9}{124\times5}{620}
\OralItem{10}{125\times2}{250}
\OralItem{11}{360\div12}{30}
\OralItem{12}{480\div16}{30}\end{oralgrid}
\par

\QHeading{口算思路补充（对应上表题号）}

\Answer{上面列出得数，下面按小题给出不同计算思路。计数单位法适用于小数，整数题可比较凑整、拆分或商不变等方法。}
\MethodTitle{第1小题：$16\times 6$}

\MethodBlock{口算思路一：按数位拆分计算}{把一个因数拆成整十、整百数与余下部分之和。}{\begin{aligned}[t]&\phantom{={}}16\times 6\\
&=6\times (10+6)\\
&=6\times 10+6\times 6\\
&=60+36\\
&=96\end{aligned}}
\MethodTitle{第6小题：$450-21$}

\MethodBlock{口算思路一：补差法}{从减数先补到一个整数，再补到被减数，两段差相加。}{\begin{aligned}[t]&\phantom{={}}450-21\\
&=(30-21)+(450-30)\\
&=9+420\\
&=429\end{aligned}}
\MethodTitle{第8小题：$95\div 5$}

\MethodBlock{口算思路一：拆分被除数}{}{\begin{aligned}[t]&\phantom{={}}95\div 5\\
&=(50+45)\div 5\\
&=50\div 5+45\div 5\\
&=10+9\\
&=19\end{aligned}}

\MethodBlock{口算思路二：利用商不变的性质}{被除数和除数同时乘同一个非零数，商不变。}{\begin{aligned}[t]&\phantom{={}}95\div 5\\
&=(95\times 2)\div (5\times 2)\\
&=190\div 10\\
&=19\end{aligned}}
\MethodTitle{第9小题：$124\times 5$}

\MethodBlock{口算思路一：按数位拆分计算}{把一个因数拆成整十、整百数与余下部分之和。}{\begin{aligned}[t]&\phantom{={}}124\times 5\\
&=5\times (100+24)\\
&=5\times 100+5\times 24\\
&=500+120\\
&=620\end{aligned}}
\MethodTitle{第10小题：$125\times 2$}

\MethodBlock{口算思路一：按数位拆分计算}{把一个因数拆成整十、整百数与余下部分之和。}{\begin{aligned}[t]&\phantom{={}}125\times 2\\
&=2\times (100+25)\\
&=2\times 100+2\times 25\\
&=200+50\\
&=250\end{aligned}}
\MethodTitle{第11小题：$360\div 12$}

\MethodBlock{口算思路一：利用商不变的性质}{被除数和除数同时乘同一个非零数，商不变。}{\begin{aligned}[t]&\phantom{={}}360\div 12\\
&=(360\times 5)\div (12\times 5)\\
&=1800\div 60\\
&=30\end{aligned}}

\MethodBlock{口算思路二：拆分除数，连续相除}{}{\begin{aligned}[t]&\phantom{={}}360\div 12\\
&=360\div (2\times 6)\\
&=360\div 2\div 6\\
&=180\div 6\\
&=30\end{aligned}}
\MethodTitle{第12小题：$480\div 16$}

\MethodBlock{口算思路一：利用商不变的性质}{被除数和除数同时乘同一个非零数，商不变。}{\begin{aligned}[t]&\phantom{={}}480\div 16\\
&=(480\times 5)\div (16\times 5)\\
&=2400\div 80\\
&=30\end{aligned}}

\MethodBlock{口算思路二：拆分除数，连续相除}{}{\begin{aligned}[t]&\phantom{={}}480\div 16\\
&=480\div (2\times 8)\\
&=480\div 2\div 8\\
&=240\div 8\\
&=30\end{aligned}}


\par\noindent

\QHeading{第2题　按要求画一画。}

\par\noindent\makebox[\linewidth][c]{
\begin{tikzpicture}[x=1bp,y=-1bp,color=black,draw=black,text=black]
\begin{scope}\clip (56.0,244.0) rectangle (525.0,367.33);
% 高清原题重排：PDF 22，原书 17。不含解答。
\hdhead{17}{相交与平行——信息窗2}
\hdsection{1}{137}{口算。}
\q{94}{184}{182}{13.5}{\(16\times6=\)}
\q{288}{184}{182}{13.5}{\(900\times0=\)}
\q{485}{184}{182}{13.5}{\(420\times20=\)}
\q{679}{184}{182}{13.5}{\(20\times300=\)}
\q{94}{228}{182}{13.5}{\(32+140=\)}
\q{288}{228}{182}{13.5}{\(450-21=\)}
\q{485}{228}{182}{13.5}{\(90\times70=\)}
\q{679}{228}{182}{13.5}{\(95\div5=\)}
\q{94}{271}{182}{13.5}{\(124\times5=\)}
\q{288}{271}{182}{13.5}{\(125\times2=\)}
\q{485}{271}{182}{13.5}{\(360\div12=\)}
\q{679}{271}{182}{13.5}{\(480\div16=\)}
\hdsection{2}{328}{按要求画一画。}
\q{94}{374}{31}{13.5}{①}
\q{127}{375}{329}{13.5}{过直线外一点，画已知直线的}
\q{127}{417}{329}{13.5}{平行线和垂线。}
\q{485}{376}{31}{13.5}{②}
\q{518}{377}{329}{13.5}{下面是一个长方形的两条相邻}
\q{518}{419}{329}{13.5}{的边，将这个长方形画完整。}
\begin{scope}[scale=.6666667,draw=black,line width=.65pt]
\fill (279,471) circle (1.6);
\draw (206,532)--(364,534);
\draw (629,463)--(629,519)--(712,519);
\draw (629,513)--(635,513)--(635,519);
\end{scope}
\hdsection{3}{554}{解决问题。}
\q{98}{600}{31}{13.5}{①}
\q{131}{601}{715}{13.5}{从王家村挖一条水渠和小河相通，怎样挖才能使水渠最短？请在图}
\q{131}{643}{715}{13.5}{中画出来。}
\q{434}{700}{112}{13.5}{王家村}
\begin{scope}[scale=.6666667,draw=black,line width=.65pt]
\fill (437,732) circle (1.5);
\draw (433,835)--(527,802)--(605,749);
\draw (440,860)--(544,824)--(619,774);
\node[rotate=19.3,inner sep=0pt,font=\normalsize] at (486,829) {小河};
\end{scope}
\q{98}{883}{31}{13.5}{②}
\q{131}{884}{716}{13.5}{点\(D\)是一处街心花园，从街心花园向\(OA\)边上修建一条长度最短的}
\q{131}{926}{716}{13.5}{小路，再过街心花园修建一条与\(OB\)边平行的小路。请在图中画出}
\q{131}{968}{716}{13.5}{这两条小路的位置。}
\q{370}{1013}{35}{13.5}{\(A\)}
\q{414}{1111}{35}{13.5}{\(O\)}
\q{576}{1111}{35}{13.5}{\(B\)}
\q{495}{1019}{35}{13.5}{\(D\)}
\begin{scope}[scale=.6666667,draw=black,line width=.65pt]
\draw (386,1015)--(426,1108)--(590,1108);
\fill (489,1031) circle (1.5);
\end{scope}

\begin{scope}[color=answerred,draw=answerred,text=answerred]\ans{134.00}{126.00}{55.33}{11}{96}
\ans{264.67}{126.00}{55.33}{11}{0}
\ans{396.00}{126.00}{55.33}{11}{8400}
\ans{524.67}{126.00}{55.33}{11}{6000}
\ans{134.00}{155.33}{55.33}{11}{172}
\ans{264.67}{155.33}{55.33}{11}{429}
\ans{396.00}{155.33}{55.33}{11}{6300}
\ans{524.67}{155.33}{55.33}{11}{19}
\ans{134.00}{183.33}{55.33}{11}{620}
\ans{264.67}{183.33}{55.33}{11}{250}
\ans{396.00}{183.33}{55.33}{11}{30}
\ans{524.67}{183.33}{55.33}{11}{30}
\draw[answerred,line width=.8pt] (133.333,313.333)--(243.333,314.726);
\draw[answerred,line width=.8pt] (186.177,300.000)--(185.367,364.000);
\draw[answerred,line width=.7pt] (190.810,355.344)--(190.878,350.011)--(185.545,349.943);
\ans{203.33}{306.00}{65.33}{9}{平行线}
\ans{186.67}{338.00}{64.67}{9}{垂线}
\draw[answerred,line width=.8pt] (419.333,308.667)--(474.667,308.667)--(474.667,346.000);
\draw[answerred,line width=.7pt] (469.333,308.667)--(469.333,314.000)--(474.667,314.000);
\draw[answerred,line width=1.1pt] (291.333,488.000)--(312.502,548.299);
\draw[answerred,line width=.7pt] (317.534,546.532)--(315.768,541.500)--(310.735,543.267);
\ans{416.67}{525.33}{131.33}{10}{答：沿村庄到河岸的\\垂线段挖，水渠最短。}
\draw[answerred,line width=1.0pt] (326.000,687.333)--(271.925,710.592);
\draw[answerred,line width=.7pt] (269.817,705.692)--(274.717,703.585)--(276.824,708.484);
\draw[answerred,line width=.8pt] (286.667,687.333)--(393.333,687.333);
\ans{262.00}{709.33}{16.67}{10}{\(H\)}
\ans{96}{766.67}{459}{10}{答：两条小路分别是\(DH\)和过\(D\)且平行于\(OB\)的直线。}
\end{scope}
\end{scope}\pgfresetboundingbox\path[use as bounding box] (56.0,244.0) rectangle (525.0,367.33);
\end{tikzpicture}}\par



\par



\par\noindent

\QHeading{第2题　补充作法}

\Answer{过指定点作平行线：可平移三角板，也可利用“两条直线都垂直于同一条直线，则它们平行”。补成长方形：可分别过两端点作已知两边的平行线，也可量出两条邻边长，在对应平行方向上截取相等长度。给定的邻边和顶点不变时，补出的长方形不因作图方法而改变。}

\QHeading{第3题　解决问题。}

\Problem{①从王家村挖一条水渠和小河相通，怎样挖才能使水渠最短？请在图中画出来。}
\par




\par\noindent\makebox[\linewidth][c]{
\begin{tikzpicture}[x=1bp,y=-1bp,color=black,draw=black,text=black]
\begin{scope}\clip (283.67,460.67) rectangle (417.67,581.34);
% 高清原题重排：PDF 22，原书 17。不含解答。
\hdhead{17}{相交与平行——信息窗2}
\hdsection{1}{137}{口算。}
\q{94}{184}{182}{13.5}{\(16\times6=\)}
\q{288}{184}{182}{13.5}{\(900\times0=\)}
\q{485}{184}{182}{13.5}{\(420\times20=\)}
\q{679}{184}{182}{13.5}{\(20\times300=\)}
\q{94}{228}{182}{13.5}{\(32+140=\)}
\q{288}{228}{182}{13.5}{\(450-21=\)}
\q{485}{228}{182}{13.5}{\(90\times70=\)}
\q{679}{228}{182}{13.5}{\(95\div5=\)}
\q{94}{271}{182}{13.5}{\(124\times5=\)}
\q{288}{271}{182}{13.5}{\(125\times2=\)}
\q{485}{271}{182}{13.5}{\(360\div12=\)}
\q{679}{271}{182}{13.5}{\(480\div16=\)}
\hdsection{2}{328}{按要求画一画。}
\q{94}{374}{31}{13.5}{①}
\q{127}{375}{329}{13.5}{过直线外一点，画已知直线的}
\q{127}{417}{329}{13.5}{平行线和垂线。}
\q{485}{376}{31}{13.5}{②}
\q{518}{377}{329}{13.5}{下面是一个长方形的两条相邻}
\q{518}{419}{329}{13.5}{的边，将这个长方形画完整。}
\begin{scope}[scale=.6666667,draw=black,line width=.65pt]
\fill (279,471) circle (1.6);
\draw (206,532)--(364,534);
\draw (629,463)--(629,519)--(712,519);
\draw (629,513)--(635,513)--(635,519);
\end{scope}
\hdsection{3}{554}{解决问题。}



\q{434}{700}{112}{13.5}{王家村}
\begin{scope}[scale=.6666667,draw=black,line width=.65pt]
\fill (437,732) circle (1.5);
\draw (433,835)--(527,802)--(605,749);
\draw (440,860)--(544,824)--(619,774);
\node[rotate=19.3,inner sep=0pt,font=\normalsize] at (486,829) {小河};
\end{scope}
\q{98}{883}{31}{13.5}{②}
\q{131}{884}{716}{13.5}{点\(D\)是一处街心花园，从街心花园向\(OA\)边上修建一条长度最短的}
\q{131}{926}{716}{13.5}{小路，再过街心花园修建一条与\(OB\)边平行的小路。请在图中画出}
\q{131}{968}{716}{13.5}{这两条小路的位置。}
\q{370}{1013}{35}{13.5}{\(A\)}
\q{414}{1111}{35}{13.5}{\(O\)}
\q{576}{1111}{35}{13.5}{\(B\)}
\q{495}{1019}{35}{13.5}{\(D\)}
\begin{scope}[scale=.6666667,draw=black,line width=.65pt]
\draw (386,1015)--(426,1108)--(590,1108);
\fill (489,1031) circle (1.5);
\end{scope}

\begin{scope}[color=answerred,draw=answerred,text=answerred]\ans{134.00}{126.00}{55.33}{11}{96}
\ans{264.67}{126.00}{55.33}{11}{0}
\ans{396.00}{126.00}{55.33}{11}{8400}
\ans{524.67}{126.00}{55.33}{11}{6000}
\ans{134.00}{155.33}{55.33}{11}{172}
\ans{264.67}{155.33}{55.33}{11}{429}
\ans{396.00}{155.33}{55.33}{11}{6300}
\ans{524.67}{155.33}{55.33}{11}{19}
\ans{134.00}{183.33}{55.33}{11}{620}
\ans{264.67}{183.33}{55.33}{11}{250}
\ans{396.00}{183.33}{55.33}{11}{30}
\ans{524.67}{183.33}{55.33}{11}{30}
\draw[answerred,line width=.8pt] (133.333,313.333)--(243.333,314.726);
\draw[answerred,line width=.8pt] (186.177,300.000)--(185.367,364.000);
\draw[answerred,line width=.7pt] (190.810,355.344)--(190.878,350.011)--(185.545,349.943);
\ans{203.33}{306.00}{65.33}{9}{平行线}
\ans{186.67}{338.00}{64.67}{9}{垂线}
\draw[answerred,line width=.8pt] (419.333,308.667)--(474.667,308.667)--(474.667,346.000);
\draw[answerred,line width=.7pt] (469.333,308.667)--(469.333,314.000)--(474.667,314.000);
\draw[answerred,line width=1.1pt] (291.333,488.000)--(312.502,548.299);
\draw[answerred,line width=.7pt] (317.534,546.532)--(315.768,541.500)--(310.735,543.267);

\draw[answerred,line width=1.0pt] (326.000,687.333)--(271.925,710.592);
\draw[answerred,line width=.7pt] (269.817,705.692)--(274.717,703.585)--(276.824,708.484);
\draw[answerred,line width=.8pt] (286.667,687.333)--(393.333,687.333);
\ans{262.00}{709.33}{16.67}{10}{\(H\)}
\ans{96}{766.67}{459}{10}{答：两条小路分别是\(DH\)和过\(D\)且平行于\(OB\)的直线。}
\end{scope}
\end{scope}\pgfresetboundingbox\path[use as bounding box] (283.67,460.67) rectangle (417.67,581.34);
\end{tikzpicture}}\par
\Answer{答：沿村庄到河岸的\par 垂线段挖，水渠最短。}
\Problem{②点\(D\)是一处街心花园，从街心花园向\(OA\)边上修建一条长度最短的小路，再过街心花园修建一条与\(OB\)边平行的小路。请在图中画出这两条小路的位置。}



\par\noindent\makebox[\linewidth][c]{
\begin{tikzpicture}[x=1bp,y=-1bp,color=black,draw=black,text=black]
\begin{scope}\clip (241.67,669.33) rectangle (412.33,770.27);
% 高清原题重排：PDF 22，原书 17。不含解答。
\hdhead{17}{相交与平行——信息窗2}
\hdsection{1}{137}{口算。}
\q{94}{184}{182}{13.5}{\(16\times6=\)}
\q{288}{184}{182}{13.5}{\(900\times0=\)}
\q{485}{184}{182}{13.5}{\(420\times20=\)}
\q{679}{184}{182}{13.5}{\(20\times300=\)}
\q{94}{228}{182}{13.5}{\(32+140=\)}
\q{288}{228}{182}{13.5}{\(450-21=\)}
\q{485}{228}{182}{13.5}{\(90\times70=\)}
\q{679}{228}{182}{13.5}{\(95\div5=\)}
\q{94}{271}{182}{13.5}{\(124\times5=\)}
\q{288}{271}{182}{13.5}{\(125\times2=\)}
\q{485}{271}{182}{13.5}{\(360\div12=\)}
\q{679}{271}{182}{13.5}{\(480\div16=\)}
\hdsection{2}{328}{按要求画一画。}
\q{94}{374}{31}{13.5}{①}
\q{127}{375}{329}{13.5}{过直线外一点，画已知直线的}
\q{127}{417}{329}{13.5}{平行线和垂线。}
\q{485}{376}{31}{13.5}{②}
\q{518}{377}{329}{13.5}{下面是一个长方形的两条相邻}
\q{518}{419}{329}{13.5}{的边，将这个长方形画完整。}
\begin{scope}[scale=.6666667,draw=black,line width=.65pt]
\fill (279,471) circle (1.6);
\draw (206,532)--(364,534);
\draw (629,463)--(629,519)--(712,519);
\draw (629,513)--(635,513)--(635,519);
\end{scope}
\hdsection{3}{554}{解决问题。}
\q{98}{600}{31}{13.5}{①}
\q{131}{601}{715}{13.5}{从王家村挖一条水渠和小河相通，怎样挖才能使水渠最短？请在图}
\q{131}{643}{715}{13.5}{中画出来。}
\q{434}{700}{112}{13.5}{王家村}
\begin{scope}[scale=.6666667,draw=black,line width=.65pt]
\fill (437,732) circle (1.5);
\draw (433,835)--(527,802)--(605,749);
\draw (440,860)--(544,824)--(619,774);
\node[rotate=19.3,inner sep=0pt,font=\normalsize] at (486,829) {小河};
\end{scope}




\q{370}{1013}{35}{13.5}{\(A\)}
\q{414}{1111}{35}{13.5}{\(O\)}
\q{576}{1111}{35}{13.5}{\(B\)}
\q{495}{1019}{35}{13.5}{\(D\)}
\begin{scope}[scale=.6666667,draw=black,line width=.65pt]
\draw (386,1015)--(426,1108)--(590,1108);
\fill (489,1031) circle (1.5);
\end{scope}

\begin{scope}[color=answerred,draw=answerred,text=answerred]\ans{134.00}{126.00}{55.33}{11}{96}
\ans{264.67}{126.00}{55.33}{11}{0}
\ans{396.00}{126.00}{55.33}{11}{8400}
\ans{524.67}{126.00}{55.33}{11}{6000}
\ans{134.00}{155.33}{55.33}{11}{172}
\ans{264.67}{155.33}{55.33}{11}{429}
\ans{396.00}{155.33}{55.33}{11}{6300}
\ans{524.67}{155.33}{55.33}{11}{19}
\ans{134.00}{183.33}{55.33}{11}{620}
\ans{264.67}{183.33}{55.33}{11}{250}
\ans{396.00}{183.33}{55.33}{11}{30}
\ans{524.67}{183.33}{55.33}{11}{30}
\draw[answerred,line width=.8pt] (133.333,313.333)--(243.333,314.726);
\draw[answerred,line width=.8pt] (186.177,300.000)--(185.367,364.000);
\draw[answerred,line width=.7pt] (190.810,355.344)--(190.878,350.011)--(185.545,349.943);
\ans{203.33}{306.00}{65.33}{9}{平行线}
\ans{186.67}{338.00}{64.67}{9}{垂线}
\draw[answerred,line width=.8pt] (419.333,308.667)--(474.667,308.667)--(474.667,346.000);
\draw[answerred,line width=.7pt] (469.333,308.667)--(469.333,314.000)--(474.667,314.000);
\draw[answerred,line width=1.1pt] (291.333,488.000)--(312.502,548.299);
\draw[answerred,line width=.7pt] (317.534,546.532)--(315.768,541.500)--(310.735,543.267);
\ans{416.67}{525.33}{131.33}{10}{答：沿村庄到河岸的\\垂线段挖，水渠最短。}
\draw[answerred,line width=1.0pt] (326.000,687.333)--(271.925,710.592);
\draw[answerred,line width=.7pt] (269.817,705.692)--(274.717,703.585)--(276.824,708.484);
\draw[answerred,line width=.8pt] (286.667,687.333)--(393.333,687.333);
\ans{262.00}{709.33}{16.67}{10}{\(H\)}

\end{scope}
\end{scope}\pgfresetboundingbox\path[use as bounding box] (241.67,669.33) rectangle (412.33,770.27);
\end{tikzpicture}}\par
\Answer{答：两条小路分别是$DH$和过$D$且平行于$OB$的直线。}


\QHeading{第3题　作图方法补充}

\Answer{方法一：用三角板直接作村庄到河岸、点D到OA所在直线的垂线段。方法二：用透明纸折出垂直方向，再把经过指定点的垂线描回。过D平行于OB的小路，可用平移三角板法，也可用两次作垂线法。这里是方法不同，不是同一个指定点能向同一条直线作多条垂线。}
